\section{The mechanical adversary}
\label{sec:algorithm}

We first specify one fixed stage of the algorithm. Its parameters are
constants independent of the play. Section~\ref{sec:uniformity} gives the
effective rule selecting an improving stage from $n$.

\subsection{Stored state}

A node $v$ contains an ordered list of positions, an interval for every free
item, fixed positions for the other items, and selected weak inequalities
$X_x\leq X_y$. Position indices are $1,\ldots,K$; positions also have
permanent identities, so insertion never changes the identity of an old
fixed position. Every position with fixed items has a child containing exactly
those items. Each original item is free at exactly one current node.

There are two kinds of record. The public order retains every actual answer.
The smaller \emph{active graph} retains selected inequalities between free
items. Comparisons involving fixed items become interval bounds. Some active
edges will be retired in favor of stronger interval bounds; their actual
comparisons are never removed from $P$.
We keep the underlying comparisons and fixed identities so their interval
consequences can be recomputed after an insertion. Separately stored
\emph{artificial bounds} are intervals imposed by the adversary itself.

An assignment is admissible if it respects these intervals, fixed values,
and inequalities. The stored conditions imply $X_x\leq X_y$ whenever
$x<_P y$ among items present in the node. Give all admissible assignments
equal weight; probabilities below always refer to this finite distribution.
The algorithm need not compute $Z_v$ globally.

Propagate lower bounds forward along inequalities and upper bounds backward.
If $[a_x,b_x]$ are the stored intervals, the resulting effective intervals are
\begin{equation}
\ell_x=\max_{y\preceq x}a_y,\qquad
r_x=\min_{y\succeq x}b_y,\qquad W_x=r_x-\ell_x+1,
\label{eq:effective}
\end{equation}
where reachability includes the item itself and fixed values are singleton
intervals. After propagation, fixed items can be eliminated from numerical
calculations. Every free item must have $W_x\geq2$.
An item newly entering a child must be publicly incomparable with all items
already there.

\subsection{The primitive move: fix in the parent, admit in the child}
\label{sec:move}

To move $x$, choose a position $c$ strictly inside its effective interval,
fix $X_x=c$, and transfer $x$ to the child at $c$ as a new free item.
Incident active edges are replaced by their interval consequences.

If the chosen position is the lower endpoint $c=\ell_x<r_x$, first insert a
new position $c^+$ immediately after it. Enlarge each artificial interval
containing $c$ to include both copies. Old fixed items stay at the old
identity $c$. Recompute comparison-induced bounds from their original
inequalities, and fix $x$ at $c^+$. At the upper endpoint use the reflected
operation, inserting immediately before it. This preserves the cost bound
for the originally selected endpoint; see Lemma~\ref{lem:split}.

A child begins with two positions. At its $j$th admission, append a top
position if necessary to give it at least $j$ positions. Leave all old
interval endpoints unchanged, and give only the incoming item the full
current list. The incoming item has no comparisons with old child items, so
this multiplies the child's count by at least $j$.

\subsection{A reproducible choice of position}
\label{sec:representatives}

For ordinary moves we will obtain an interval $I$ of $m$ positions such that
every $t\in I$ has probability at least $1/(16m)$. Certain repairs use a
different constant $V$ in place of $16$. In either case the following rule
chooses a position that can be reused.

Maintain binary-counter batches of positions: a position is born in a
singleton batch, and equal-size batches merge until there is at most one of
each size. A batch's members keep their relative position order for its
lifetime, even if other positions are inserted between them.
Within a batch, the aligned block of size $2^k$ with indices
$j2^k+1,\ldots,(j+1)2^k$ has representative its sole member if $k=0$,
and its member at index $j2^k+2^{k-1}$ otherwise.

For $m\geq8$, delete the first and last positions of $I$. Intersect the
remaining interval with every current batch, decompose each intersection
into maximal aligned blocks, and choose a largest block. Fix at its
representative. For $m<8$, use the first position of $I$ and perform an
endpoint insertion if needed. Record this exceptional choice as block size
one. Initially present positions are entered into batches in their order.

A scale group below uses a restricted version. During that group a position
cannot be chosen twice as a singleton-block representative, or twice as a
larger-block representative. On later choices discard earlier chosen blocks
of the same kind before taking a largest block. A position represents at
most one larger block in a fixed batch. The width separation below ensures
enough blocks remain. The group's parent list and batches stay unchanged.

\subsection{Stage constants and local queries}
\label{sec:constants}

Choose integers $D,H\geq4$ and a rational $0<\eta\leq1/16$. Set
\begin{align}
d&=D+1,&h&=2H,& M&=(d+1)^{h+1},&
B&=2M,& C_*&=2\cdot4^B .
\label{eq:base-parameters}
\end{align}
Choose integers $L\geq4$, $J\geq4$ with $C_*/L\leq\eta$, and set
\begin{equation}
A=128J(L+1)^{2h+3},\qquad R=A^{J+1}.
\label{eq:scale-parameters}
\end{equation}
For these $d,h,R,\eta$, compute the radius $\rho$ and width threshold $W_0$
from Appendix~\ref{app:local}. Increase $W_0$ to at least $32$ if needed.
Put $G=4^B W_0$.

After repairs, active degree is less than $D$ and active height (vertices
on a longest directed path) is less than $H$. During repairs we may use
the larger bounds $d,h$ above. Each item then has at most $B$ other active
predecessors and successors.

A \textsc{LocalQuery} on one or two items retains their radius-$\rho$
neighborhood in the active graph, together with the \emph{globally
propagated} intervals \eqref{eq:effective}. It counts this finite restricted
system exactly using the subset dynamic program in
Section~\ref{sec:counter}. It can return either a weak comparison probability
or a coordinate cumulative probability. After width repairs these estimates
have error at most $\eta$. An isolated item is an independent uniform
coordinate and is handled exactly, at any width.

An \textsc{OrdinaryMove} finds the first crossings $a,b$ of $1/4,3/4$
in one such local cumulative distribution, applies the representative rule
to $[a,b]$, and moves the item. Since $\eta\leq1/16$, every position in
that interval has true probability at least $1/(16(b-a+1))$.

\subsection{The repair procedure}
\label{sec:repairs}

The following rules act only in the node that has just answered a comparison.
Recompute effective intervals after every change.

\paragraph{Width repairs.}
Repeatedly take the first applicable rule, ignoring isolated items.
\begin{enumerate}[label=\textbf{W\arabic*.},leftmargin=11mm]
\item \textbf{Narrow item.}
If $W_x<W_0$, let
\[
 I_x=\left[\ell_x+\left\lfloor\frac{W_x-1}{4}\right\rfloor,\,
 r_x-\left\lfloor\frac{W_x-1}{4}\right\rfloor\right].
\]
Choose its first position, make it interior if needed, and move $x$.
This is a \emph{cheap move}; its parent retains at least $1/G$.

\item \textbf{An edge across unequal widths.}
For an edge $u\to v$ with $W_u>LW_v$, impose the artificial bounds
$X_u\leq\ell_v-1$ and $X_v\geq\ell_v$, then retire the active edge.
If $W_v>LW_u$, instead impose $X_u\leq r_u$ and $X_v\geq r_u+1$.
These strict cuts preserve the answered comparison and retain at least
$1-\eta$ of the count.

\item \textbf{A component across too many scales.}
If a component's largest width exceeds $R$ times its smallest width $w$,
take a shortest path from a smallest-width item to a largest-width item.
For $i=0,\ldots,J-1$, select the first item of width at least $A^i w$.
Process the selected items in increasing width, choosing representatives
from their central intervals $I_x$ by the restricted group rule.
These are \emph{scale moves}. They use interior positions, insert no
parent positions, and send at most two items to each child.
\end{enumerate}
Return to W1 after each rule or completed group. The selected scale items
retain their intervals and incident edges until processed.

\paragraph{Degree and path repairs.}
First finish all width repairs.
If an item currently has degree at least $D$, perform an ordinary move of
that item, repair widths, and retest degrees. Otherwise, find a longest
directed path by topological dynamic programming. If it contains at least
$H$ free items, reserve its first $H$ items. Process them in path
order: before each ordinary move repair widths; skip reserved items already
moved by those repairs. Move a reserved item even if it has since become
isolated. Complete this reservation before choosing a new path.
Repeat until degree is less than $D$ and height is less than $H$.

No child is repaired merely because it receives an item: the new item is
isolated there and its admission changes none of the old distributions.

\subsection{Answering an actual comparison}
\label{sec:answer}

An answer at a node is \emph{safe} if it neither creates a new public
relation between two old items fixed at the same position, nor reduces any
free effective interval to a singleton, and its weak constraints remain
feasible. Test this by public transitive closure and interval propagation
under each proposed answer.

\begin{algorithm}[ht]
\caption{\textsc{Answer}$(x,y)$}
\label{alg:answer}
\begin{algorithmic}[1]
\If{$P$ already implies the comparison}
  \State \Return its implied answer.
\EndIf
\State Set $v$ to the root.
\While{$x,y$ are both fixed at the same position $c$ in $v$}
  \State Set $v$ to the child at $c$.
\EndWhile
\State Test safety of $x<y$ and $y<x$ at $v$.
\If{one answer is unsafe}
  \State Choose the safe answer retaining every old assignment.
\Else
  \State Estimate $\Prob(X_x\leq X_y)$ and $\Prob(X_y\leq X_x)$ locally.
  \State Choose the answer with larger estimate.
\EndIf
\State Add the chosen strict answer to $P$ and its weak constraint at $v$.
\State Run width, degree, and path repairs at $v$.
\State \Return the chosen answer.
\end{algorithmic}
\end{algorithm}

Only a comparison with two free endpoints adds an active edge; a fixed
endpoint supplies a bound. Routing through ancestors changes none of their
counts. Lemma~\ref{lem:routing} proves that the safe all-retaining branch
always exists when needed.

The root starts with $n$ positions, all $n$ items free on the full interval,
and no edges. For $n\leq1$ there is nothing to answer. If no stage fits the
size caps of Section~\ref{sec:uniformity}, use a fixed label order.
After $2n\lceil\logtwo(n+1)\rceil$ actual queries, the algorithm may fix any
total order extending $P$ and answer according to it thereafter. This final
rule bounds storage even against a sorter making indefinitely many redundant
queries; by then the desired comparison threshold has already been exceeded.
